\documentclass[10pt,a4paper]{amsart}
\usepackage[margin=19mm]{geometry}
\usepackage{amsmath,amssymb,mathtools}
\usepackage[scr=rsfs,cal=euler]{mathalfa}
\usepackage{xfrac}
\usepackage[unicode,hidelinks,bookmarks=false]{hyperref}
\newcommand{\ie}{i.e.\ }
\newcommand{\cf}{cf.\ }
\newcommand{\resp}{respectively\ }
\newcommand{\Subsection}{\subsection}
\providecommand{\nobreakdash}{\nobreak\hbox{-}}
\setlength{\emergencystretch}{2em}
\multlinegap0pt
\usepackage{amssymb}
\usepackage{tikz}
\newtheorem{lem}{Lemma}
\title{Hilbert modular Eisenstein congruences of local origin.}
\author{Dan Fretwell, Jenny Roberts}
\date{Source: arXiv:2411.06987v2 (CC BY 4.0)}
\begin{document}
\maketitle

\noindent\emph{Source and licence.} Hilbert modular Eisenstein congruences of local origin.. Version arXiv:2411.06987v2 (CC BY 4.0), distributed by arXiv under the Creative Commons Attribution 4.0 International licence, http://creativecommons.org/licenses/by/4.0/, as recorded in the arXiv licence metadata for this exact version and shown on the abstract page https://arxiv.org/abs/2411.06987v2. Full licence notice: \texttt{licenses/arxiv-2411.06987-CC-BY-4.0.txt}.\par\smallskip

\noindent\textit{Editorial scope.} Counted unit: the article's lem lem (source0.tex lines 348--358). The article's complete original proof of it is delivered with it (source0.tex lines 359--420, one \texttt{proof} environment of 62 lines). No mathematics is written, repaired or re-derived: the statement and the proof are the article's own lines, in the article's own notation, and no retained line is substituted or re-flowed.  Retained uncounted as the source-local context the proof needs: lines 320--323.  Nothing was omitted from any retained span, so the package carries no omission record.  No retained line contains a citation, so the wrapper carries no bibliography and no bibliography entry.  No retained line contains a figure environment, an image-inclusion command or drawing code, so no image is delivered and the package declares no asset dependency.  Editorial formatting only, in addition: an AMS \texttt{amsart} preamble that restates the article's own declarations and macros used by the retained lines, namely the environments \texttt{lem} on separate counters, together with the standard packages those lines need.  The retained environments print in this document as Lemma 1 in order of appearance, whereas the article prints the same items with the numbering of its own full document.  The complete native source of the article as distributed in the arXiv e-print for this exact version is bundled unchanged as \texttt{sources/168334/source0.tex}.
\begin{equation}
\label{eq:Edelta}
    E^\delta = E_k^{(\mathcal{O})}(\eta, \psi) - \delta E_k^{(\mathfrak{p})}(\eta, \psi) \in E_{\textbf{k}}^{\text{new},(\mathfrak{p})}(\mathfrak{m},\phi)\subseteq E_{\textbf{k}}(\mathfrak{mp},\phi).
\end{equation}

\begin{lem}
    The Eisenstein series $E^\delta$ for $\delta \in \{\eta(\mathfrak{p}), \psi(\mathfrak{p})N(\mathfrak{p})^{k-1}\} $ are normalised eigenforms with eigenvalues given by:
    \begin{equation*}
        T_{\mathfrak{q}}(E^\delta) = \begin{cases}
            (\eta(\mathfrak{q})+\psi(\mathfrak{q})N(\mathfrak{q})^{k-1}) E^\delta \ &\text{if } \mathfrak{q} \neq \mathfrak{p}\\
            \varepsilon E^\delta \ &\text{if } \mathfrak{q} = \mathfrak{p}
        \end{cases}
    \end{equation*}
    where $\varepsilon = \eta(\mathfrak{p})+\psi(\mathfrak{p})N(\mathfrak{p})^{k-1} - \delta$.
    
\end{lem}

\begin{proof}
From Equation $\eqref{eq:Edelta}$, we have that the $\mathfrak{q}^{\text{th}}$ Fourier coefficient of $E^\delta$, for a prime ideal $\mathfrak{q}$, is given by:
\begin{equation}
    c(\mathfrak{q}, E^\delta) = \begin{cases}
    \eta(\mathfrak{q})+\psi(\mathfrak{q})N(\mathfrak{q})^{k-1}\ &\text{if } \mathfrak{q} \neq \mathfrak{p} \\
            \varepsilon \ &\text{if } \mathfrak{q} = \mathfrak{p}
    \end{cases}
\end{equation}
    and so we just need to check that $E^\delta$ is an eigenform. To do so, we first rewrite $E^\delta$ in the following form:
    \begin{equation}
    \label{eq:HeckeE}
        E^\delta = (T_\mathfrak{p} - \delta)E^{(\mathfrak{p})}_k(\eta, \psi).
    \end{equation}
    Then, we show that the eigenvalues of the right hand side of Equation \eqref{eq:HeckeE} match those in the lemma.

    We will split this into cases, evaluating $T_\mathfrak{q}E^\delta$ for $\mathfrak{q} \neq \mathfrak{p}$ and $T_\mathfrak{p}E^\delta$. 

    \vspace{0.2in}
    \textbf{Case $1$: $\mathfrak{q} \neq \mathfrak{p}$}
    
    \vspace{0.1in}
    Note that the Hecke operators $T_\mathfrak{q}$ and $T_\mathfrak{p}$ commute since $\mathfrak{q} \neq \mathfrak{p}$, so we have 
    \begin{align*}
        T_\mathfrak{q}E^\delta &= T_\mathfrak{q}(T_\mathfrak{p} - \delta)E^{(\mathfrak{p})}_k(\eta, \psi) \\
        &= (T_\mathfrak{p} - \delta)T_\mathfrak{q}E^{(\mathfrak{p})}_k(\eta, \psi).
    \end{align*}
     Now we just need to determine $T_\mathfrak{q}E^{(\mathfrak{p})}_k(\eta, \psi)$. Below, we write $\tilde{\phi}$ for the lift of character $\phi$ to modulus $\mathfrak{nq}$.

     \begin{align*}
     c(\mathfrak{n},T_\mathfrak{q}E^{(\mathfrak{p})}_k(\eta, \psi)) &= c(\mathfrak{nq}, E^{(\mathfrak{p})}_k(\eta, \psi)) + \tilde{\phi}(\mathfrak{q})c(\mathfrak{nq}^{-1}, E^{(\mathfrak{p})}_k(\eta, \psi)) \\
         &= c(\mathfrak{nqp}^{-1}, E_k(\eta,\psi)) + \phi(\mathfrak{q})c(\mathfrak{nq}^{-1}\mathfrak{p}^{-1}, E_k(\eta, \psi)) \\
         &= c(\mathfrak{np}^{-1}, T_\mathfrak{q}E_k(\eta,\psi)) \\
         &= (\eta(\mathfrak{q})+\psi(\mathfrak{q})N(\mathfrak{q})^{k-1})c(\mathfrak{np}^{-1}, E_k(\eta,\psi)) \\
         &=(\eta(\mathfrak{q})+\psi(\mathfrak{q})N(\mathfrak{q})^{k-1})c(\mathfrak{n}, E_k^{(\mathfrak{p})}(\eta,\psi)).
    \end{align*}

\textbf{Case $2$: $\mathfrak{q} = \mathfrak{p}$}

\vspace{0.1in}
In this case we have
\begin{equation*}
T_\mathfrak{p}E^\delta = (T_\mathfrak{p}^2-\delta T_\mathfrak{p})E_k^{(\mathfrak{p})}(\eta, \psi),
\end{equation*}
and so we will first determine $T_\mathfrak{p}E_k^{(\mathfrak{p})}(\eta, \psi)$. Note that
\begin{align*}
    c(\mathfrak{n}, T_\mathfrak{p}E_k^{(\mathfrak{p})}(\eta, \psi)) &= c(\mathfrak{np}, E_k^{(\mathfrak{p})}(\eta, \psi)) \\
    &= c(\mathfrak{n}, E_k(\eta, \psi)),
\end{align*}
hence we have that
\begin{align*}
   c(\mathfrak{n}, T_\mathfrak{p}^2E_k^{(\mathfrak{p})}(\eta, \psi)) &= c(\mathfrak{n}, T_\mathfrak{p}E_k(\eta, \psi)) \\
   &=(\eta(\mathfrak{p}) + \psi(\mathfrak{p})N(\mathfrak{p})^{k-1})c(\mathfrak{n}, E_k(\eta, \psi)).
\end{align*}
Consequently,
\begin{align*}
    c(\mathfrak{n},T_\mathfrak{p}E^\delta) &= c(\mathfrak{n}, (T_\mathfrak{p}^2-\delta T_\mathfrak{p})E_k^{(\mathfrak{p})}(\eta, \psi)) \\
    &= (\eta(\mathfrak{p}) + \psi(\mathfrak{p})N(\mathfrak{p})^{k-1})c(\mathfrak{n}, E_k(\eta, \psi)) -\phi(\mathfrak{p})N(\mathfrak{p})^{k-1}c(\mathfrak{np}^{-1}, E_k(\eta, \psi)) - \delta c(\mathfrak{n}, E_k(\eta, \psi)) \\
    &=(\eta(\mathfrak{p}) + \psi(\mathfrak{p})N(\mathfrak{p})^{k-1} - \delta)c(\mathfrak{n}, T_{\mathfrak{p}}E_k^{(\mathfrak{p})}(\eta, \psi)) - \phi(\mathfrak{p})N(\mathfrak{p})^{k-1}c(\mathfrak{n}, E_k^{(\mathfrak{p})}(\eta, \psi))\\
    &= \varepsilon c(\mathfrak{n}, (T_\mathfrak{p} -\delta ) E_k^{(\mathfrak{p})}(\eta, \psi)) \\
    &= \varepsilon c(\mathfrak{n}, E^\delta).
\end{align*}
\end{proof}

\end{document}
